\subsection{Sprint 3}

Na Sprint 3, a prioridade passou a ser manter o ritmo do projeto e evitar que as tarefas atrasassem. Depois da reorganização em \textit{squads}, passei a observar mais de perto como os AGES III, II e I estavam se saindo, tentando manter a comunicação ativa e verificar se aquilo que estávamos fazendo ainda estava de acordo com o que todos esperavam. Acredito que essa sprint teve um papel mais forte de gestão do time do que de definição inicial de escopo.

Durante esse período, consegui gerenciar os times e manter em ordem aquilo que estava sendo feito. Revisei as tarefas abertas, observei o funcionamento das \textit{squads} e procurei garantir que a forma de organização continuasse fazendo sentido para o grupo. Também fiquei atento para que as entregas não se afastassem das necessidades do projeto e dos alinhamentos anteriores, porque nessa etapa já era importante evitar mudanças de direção que pudessem comprometer a reta final.

Mesmo com esse acompanhamento, no final da sprint uma tarefa ficou pendente, pois a pessoa na qual tinha se comprometido em fazê-la sumiu e não deu respostas. Como ela precisava ser resolvida para que o projeto não ficasse parado, acabei assumindo o papel de desenvolvedor e fazendo essa tarefa. Não era o cenário ideal, porque o objetivo era que cada frente conseguisse caminhar com autonomia e aprendesse com as tarefas, mas naquele momento entendi que era mais importante destravar a entrega e impedir que o problema se acumulasse para a etapa final.

A dificuldade mais visível foi a falta de interesse de alguns colegas. Esse comportamento já vinha aparecendo em outros momentos, mas na Sprint 3 ficou mais evidente porque estávamos nos aproximando da reta final e qualquer atraso passava a pesar mais. Foi frustrante perceber que, mesmo com organização e cobrança, nem todos mantinham o mesmo nível de compromisso com as tarefas. Isso exigiu mais atenção da minha parte para acompanhar o andamento e identificar rapidamente o que poderia virar bloqueio.

Em termos de aprendizado, essa sprint não trouxe uma lição nova tão forte quanto as anteriores. Foi mais uma etapa de manutenção, gestão e reação aos problemas que surgiram. Para a sprint final, ficou como tarefa verificar o deploy, observar como o time estava se saindo e dar um gás para terminar o projeto. Também assumi para mim que, se fosse necessário codar para resolver algum problema ou evitar atraso, eu colocaria a mão na massa e desenvolveria.
